\documentclass[a4paper,12pt]{article}

% -------------------------------------------------
% Кодировка, язык, математика
% -------------------------------------------------
\usepackage[T2A]{fontenc}
\usepackage[utf8]{inputenc}
\usepackage[russian]{babel}

\usepackage{amsmath,amssymb,amsthm,mathtools}
\usepackage{icomma}

% -------------------------------------------------
% Типографика
% -------------------------------------------------
\usepackage{helvet}
\renewcommand{\familydefault}{\sfdefault}

\usepackage[a4paper,margin=20mm,top=19mm,bottom=20mm]{geometry}
\usepackage{microtype}
\usepackage{parskip}
\usepackage{setspace}
\onehalfspacing

\usepackage{xcolor}
\definecolor{accent}{HTML}{008080}
\definecolor{darkaccent}{HTML}{004D4D}
\definecolor{textgray}{HTML}{333333}
\definecolor{softaccent}{HTML}{EAF4F4}
\definecolor{softgray}{HTML}{F5F6F6}
\definecolor{warning}{HTML}{8A5B38}

\color{textgray}

% -------------------------------------------------
% Графика
% -------------------------------------------------
\usepackage{tikz}
\usetikzlibrary{calc}
\usetikzlibrary{arrows.meta,positioning,shapes.geometric}

\usepackage{tkz-euclide}

\usepackage{pgfplots}
\pgfplotsset{compat=1.18}

% -------------------------------------------------
% Таблицы, списки, оформление
% -------------------------------------------------
\usepackage{tabularx,booktabs,longtable}
\usepackage{enumitem}
\usepackage[most]{tcolorbox}

\usepackage{hyperref}
\hypersetup{
    colorlinks=true,
    linkcolor=darkaccent,
    urlcolor=darkaccent
}

\usepackage{fancyhdr}
\pagestyle{fancy}
\fancyhf{}
\fancyhead[L]{\small Векторы на координатной плоскости}
\fancyhead[R]{\small Анастасия Павлова}
\fancyfoot[C]{\small\thepage}
\renewcommand{\headrulewidth}{0.3pt}
\renewcommand{\footrulewidth}{0pt}

% -------------------------------------------------
% Служебные команды
% -------------------------------------------------
\newcommand{\vect}[1]{\overrightarrow{#1}}
\newcommand{\va}{\vec a}
\newcommand{\vb}{\vec b}
\newcommand{\vc}{\vec c}
\newcommand{\vd}{\vec d}

\newcommand{\tg}{\mathop{\mathrm{tg}}\nolimits}
\newcommand{\ctg}{\mathop{\mathrm{ctg}}\nolimits}
\newcommand{\arctg}{\mathop{\mathrm{arctg}}\nolimits}

\newcommand{\sectionline}{%
    \par\noindent
    {\color{accent}\rule{\linewidth}{0.6pt}}
    \par
}

% -------------------------------------------------
% Лёгкие учебные блоки
% -------------------------------------------------
\tcbset{
    enhanced,
    breakable,
    colback=white,
    colframe=accent,
    boxrule=0.55pt,
    arc=1mm,
    left=3mm,
    right=3mm,
    top=2mm,
    bottom=2mm,
    before skip=5pt,
    after skip=6pt
}

\newtcolorbox{keybox}[1][]{
    colback=softaccent,
    colframe=accent,
    title={#1},
    fonttitle=\bfseries,
    coltitle=darkaccent
}

\newtcolorbox{examplebox}[1][]{
    colback=white,
    colframe=darkaccent,
    title={#1},
    fonttitle=\bfseries,
    coltitle=darkaccent
}

\newtcolorbox{warningbox}[1][]{
    colback=softgray,
    colframe=warning,
    title={#1},
    fonttitle=\bfseries,
    coltitle=warning
}

% -------------------------------------------------
% Стили блок-схем
% -------------------------------------------------
\tikzset{
    flowstart/.style={
        ellipse,
        draw=darkaccent,
        line width=0.8pt,
        minimum width=44mm,
        minimum height=12mm,
        align=center,
        fill=softaccent,
        font=\small
    },
    flowaction/.style={
        rectangle,
        rounded corners=2mm,
        draw=accent,
        line width=0.7pt,
        minimum width=78mm,
        text width=68mm,
        minimum height=14mm,
        align=center,
        fill=white,
        font=\small
    },
    flowcondition/.style={
        diamond,
        aspect=2.2,
        draw=accent,
        line width=0.7pt,
        text width=38mm,
        align=center,
        inner sep=1.5mm,
        fill=softgray,
        font=\small
    },
    flowarrow/.style={
        -{Stealth[length=3mm,width=2mm]},
        line width=0.8pt,
        draw=darkaccent
    }
}

% -------------------------------------------------
% Документ
% -------------------------------------------------
\begin{document}

% =================================================
% ТИТУЛЬНЫЙ ЛИСТ
% =================================================
\begin{titlepage}
\thispagestyle{empty}

\begin{center}

\vspace*{22mm}

{\color{darkaccent}\Large\bfseries ЧЕК-ЛИСТ}

\vspace{9mm}

{\Huge\bfseries Векторы на координатной плоскости}

\vspace{5mm}

{\Large Координаты, длина, действия с векторами\\
и скалярное произведение}

\vspace{14mm}

{\color{accent}\rule{0.70\linewidth}{0.8pt}}

\vspace{14mm}

{\large
\textbf{Формат:} подготовка к ЕГЭ
}

\vspace{5mm}

{\large
\textbf{Ученица:} Анастасия Павлова
}

\vspace{5mm}

{\large
\textbf{Дата:} \today
}

\vfill

{\large
Лёвин Артём Александрович\\[2mm]
\textbf{эксклюзивно для Анастасии Павловой}
}

\vspace{14mm}

\begin{tikzpicture}[remember picture,overlay]
    \draw[
        accent,
        line width=1.15pt,
        opacity=0.60
    ]
    ($(current page.south west)+(18mm,17mm)$)
    .. controls
    ($(current page.south west)+(68mm,29mm)$)
    and
    ($(current page.south east)+(-68mm,7mm)$)
    ..
    ($(current page.south east)+(-18mm,20mm)$);
\end{tikzpicture}

\end{center}
\end{titlepage}

% =================================================
% ОСНОВНОЙ ЧЕК-ЛИСТ
% =================================================

\section*{\color{darkaccent}1. Координаты вектора}

\sectionline

Пусть
\[
A(x_A;y_A), \qquad B(x_B;y_B).
\]

Тогда координаты вектора, направленного из точки \(A\) в точку \(B\),
вычисляются по формуле
\[
\boxed{
\vect{AB}=
\bigl(x_B-x_A\,;\,y_B-y_A\bigr).
}
\]

\begin{keybox}[Главная интерпретация]
Координаты вектора показывают его \textbf{смещение}:

\[
\begin{array}{ccl}
x>0 &\Longrightarrow& \text{смещение вправо},\\
x<0 &\Longrightarrow& \text{смещение влево},\\
y>0 &\Longrightarrow& \text{смещение вверх},\\
y<0 &\Longrightarrow& \text{смещение вниз}.
\end{array}
\]
\end{keybox}

Например, если вектор перемещает точку на \(6\) единиц вправо
и на \(1\) единицу вверх, то
\[
\vec a=(6;1).
\]

Если направление вектора развернуть, обе координаты меняют знак:
\[
\vect{BA}=-\vect{AB}.
\]

\begin{center}
\begin{tikzpicture}[scale=0.78,>=Stealth]
    \draw[step=1cm,very thin,gray!25] (-1,-1) grid (8,5);
    \draw[->,line width=0.9pt] (-1,0)--(8.3,0) node[right] {$x$};
    \draw[->,line width=0.9pt] (0,-1)--(0,5.3) node[above] {$y$};

    \foreach \x in {1,...,7}
        \draw (\x,0.08)--(\x,-0.08) node[below=2pt] {\small \x};

    \foreach \y in {1,...,4}
        \draw (0.08,\y)--(-0.08,\y) node[left=2pt] {\small \y};

    \coordinate (A) at (1,1);
    \coordinate (B) at (7,3);

    \fill[darkaccent] (A) circle (2.1pt);
    \fill[darkaccent] (B) circle (2.1pt);

    \draw[->,accent,line width=1.5pt] (A)--(B);

    \node[below left] at (A) {$A(1;1)$};
    \node[above right] at (B) {$B(7;3)$};

    \draw[dashed,gray] (A)--(7,1);
    \draw[dashed,gray] (7,1)--(B);

    \node[below] at (4,1) {\small \(+6\)};
    \node[right] at (7,2) {\small \(+2\)};
\end{tikzpicture}
\end{center}

Здесь
\[
\vect{AB}=(7-1;3-1)=(6;2).
\]

\begin{warningbox}[Типичная ошибка]
Порядок вычитания принципиален:
\[
\boxed{\text{конец} - \text{начало}.}
\]

Для \(\vect{AB}\) сначала используются координаты точки \(B\),
затем координаты точки \(A\).
\end{warningbox}

% -------------------------------------------------

\section*{\color{darkaccent}2. Длина вектора}

\sectionline

Если
\[
\vec a=(x;y),
\]
то его длина равна
\[
\boxed{
|\vec a|=\sqrt{x^2+y^2}.
}
\]

Это формула расстояния, полученная из теоремы Пифагора.

\begin{examplebox}[Пример]
Пусть
\[
\vec a=(3;4).
\]
Тогда
\[
|\vec a|
=
\sqrt{3^2+4^2}
=
\sqrt{9+16}
=
5.
\]
\end{examplebox}

Если даны две точки, сначала найдите координаты вектора.

\begin{examplebox}[Длина вектора по координатам его концов]
Пусть
\[
M(3;4),\qquad T(-3;-5).
\]

Тогда
\[
\vect{MT}
=
(-3-3;-5-4)
=
(-6;-9).
\]

Следовательно,
\[
|\vect{MT}|
=
\sqrt{(-6)^2+(-9)^2}
=
\sqrt{36+81}
=
\sqrt{117}
=
3\sqrt{13}.
\]
\end{examplebox}

\begin{keybox}[Квадрат длины]
Иногда в задаче требуется найти именно
\[
|\vec a|^2.
\]

Тогда извлекать корень вообще не требуется:
\[
\boxed{
|\vec a|^2=x^2+y^2.
}
\]

Например, для
\[
\vec a=(-6;2)
\]
получаем
\[
|\vec a|^2=(-6)^2+2^2=40.
\]
\end{keybox}

% -------------------------------------------------

\section*{\color{darkaccent}3. Осторожно с квадратными корнями}

\sectionline

В задачах на длину вектора часто появляются квадратные корни.

Допустимо:
\[
\sqrt{ab}=\sqrt a\sqrt b,
\qquad a\ge0,\quad b\ge0.
\]

Например,
\[
\sqrt{75}
=
\sqrt{25\cdot3}
=
5\sqrt3.
\]

Также
\[
\sqrt{\frac ab}
=
\frac{\sqrt a}{\sqrt b},
\qquad
a\ge0,\quad b>0.
\]

\begin{warningbox}[Так делать нельзя]
В общем случае
\[
\boxed{
\sqrt{a+b}\ne\sqrt a+\sqrt b
}
\]
и
\[
\boxed{
\sqrt{a-b}\ne\sqrt a-\sqrt b.
}
\]

Например,
\[
\sqrt{75}
=
\sqrt{100-25}
\ne
10-5.
\]

Кроме того,
\[
\boxed{\sqrt{x^2}=|x|.}
\]
\end{warningbox}

% -------------------------------------------------

\section*{\color{darkaccent}4. Сложение и вычитание векторов}

\sectionline

Пусть
\[
\vec a=(x_1;y_1),
\qquad
\vec b=(x_2;y_2).
\]

Тогда
\[
\boxed{
\vec a+\vec b=
(x_1+x_2\,;\,y_1+y_2)
}
\]

и
\[
\boxed{
\vec a-\vec b=
(x_1-x_2\,;\,y_1-y_2).
}
\]

\begin{examplebox}[Пример]
Для
\[
\vec a=(1;2),
\qquad
\vec b=(3;4)
\]
получаем
\[
\vec a+\vec b
=
(1+3;2+4)
=
(4;6),
\]
а
\[
\vec a-\vec b
=
(1-3;2-4)
=
(-2;-2).
\]
\end{examplebox}

Можно складывать любое конечное число векторов:
\[
\vec a+\vec b+\vec c.
\]

Координаты по-прежнему складываются отдельно:
\[
(x_a+x_b+x_c\,;\,y_a+y_b+y_c).
\]

\begin{keybox}[Стратегия]
Сначала найдите координаты каждого отдельного вектора.

После этого работайте с координатами как с обычными числовыми парами:
складывайте, вычитайте и подставляйте в нужную формулу.
\end{keybox}

% -------------------------------------------------

\section*{\color{darkaccent}5. Умножение вектора на число}

\sectionline

Если
\[
\vec a=(x;y),
\]
то
\[
\boxed{
k\vec a=(kx;ky).
}
\]

Например,
\[
\vec a=(1;2)
\]
и
\[
3\vec a=(3;6).
\]

Если коэффициент отрицательный,
\[
-3\vec a=(-3;-6).
\]

\begin{keybox}[Что происходит геометрически]
При умножении на число меняется длина вектора.

Если коэффициент отрицательный, дополнительно меняется направление
на противоположное.
\end{keybox}

% -------------------------------------------------

\section*{\color{darkaccent}6. Скалярное произведение}

\sectionline

Для векторов
\[
\vec a=(x_1;y_1),
\qquad
\vec b=(x_2;y_2)
\]
скалярное произведение можно вычислить через координаты:
\[
\boxed{
\vec a\cdot\vec b=x_1x_2+y_1y_2.
}
\]

Скалярное произведение --- это \textbf{число}.

\begin{examplebox}[Пример]
Пусть
\[
\vec p=(1;2),
\qquad
\vec q=(3;4).
\]

Тогда
\[
\vec p\cdot\vec q
=
1\cdot3+2\cdot4
=
11.
\]
\end{examplebox}

Вторая важная формула:
\[
\boxed{
\vec a\cdot\vec b
=
|\vec a|\,|\vec b|\cos\varphi,
}
\]
где \(\varphi\) --- угол между ненулевыми векторами.

Отсюда
\[
\boxed{
\cos\varphi
=
\frac{\vec a\cdot\vec b}
{|\vec a|\,|\vec b|}.
}
\]

Эта формула применима, если
\[
|\vec a|\ne0
\qquad\text{и}\qquad
|\vec b|\ne0.
\]

Для нулевого вектора угол с другим вектором в данном смысле
не определяется.



% -------------------------------------------------

\section*{\color{darkaccent}7. Косинус угла между векторами}

\sectionline

Рассмотрим
\[
\vec k=(-1;2),
\qquad
\vec t=(-2;3).
\]

Сначала вычислим скалярное произведение:
\[
\vec k\cdot\vec t
=
(-1)(-2)+2\cdot3
=
8.
\]

Длины:
\[
|\vec k|
=
\sqrt{(-1)^2+2^2}
=
\sqrt5,
\]
\[
|\vec t|
=
\sqrt{(-2)^2+3^2}
=
\sqrt{13}.
\]

Следовательно,
\[
\cos\varphi
=
\frac8{\sqrt5\sqrt{13}}
=
\frac8{\sqrt{65}}.
\]

При необходимости знаменатель можно рационализировать:
\[
\frac8{\sqrt{65}}
=
\frac{8\sqrt{65}}{65}.
\]

\begin{keybox}[Полезный частный случай]
Если для двух ненулевых векторов
\[
\vec a\cdot\vec b=0,
\]
то
\[
\cos\varphi=0.
\]

Следовательно,
\[
\varphi=90^\circ,
\]
то есть векторы перпендикулярны.
\end{keybox}

И наоборот, для ненулевых векторов:
\[
\boxed{
\vec a\perp\vec b
\quad\Longleftrightarrow\quad
\vec a\cdot\vec b=0.
}
\]

% -------------------------------------------------

\section*{\color{darkaccent}8. Составные выражения с векторами}

\sectionline

В экзаменационной задаче может встретиться выражение вида
\[
(\vec a-\vec b)\cdot\vec c.
\]

Работайте по уровням.

\begin{enumerate}[label=\arabic*.,leftmargin=8mm,itemsep=2pt]
    \item Найдите координаты \(\vec a,\vec b,\vec c\).
    \item Вычислите координаты \(\vec a-\vec b\).
    \item Выполните скалярное произведение полученного вектора с \(\vec c\).
\end{enumerate}

\begin{examplebox}[Пример]
Пусть
\[
\vec a=(4;6),
\qquad
\vec b=(6;-2),
\qquad
\vec c=(1;-4).
\]

Тогда
\[
\vec a-\vec b
=
(4-6;6-(-2))
=
(-2;8).
\]

Следовательно,
\[
(\vec a-\vec b)\cdot\vec c
=
(-2)\cdot1+8\cdot(-4)
=
-34.
\]
\end{examplebox}

\begin{warningbox}[Следите за знаками]
Особенно опасны выражения вида
\[
y_1-y_2
\]
при отрицательном \(y_2\).

Например,
\[
6-(-2)=8.
\]

Скобки вокруг отрицательных координат лучше записывать всегда.
\end{warningbox}

% -------------------------------------------------

\section*{\color{darkaccent}9. Быстрый чек-лист перед ответом}

\sectionline

\begin{enumerate}[label=\(\square\),leftmargin=8mm,itemsep=4pt]
    \item Для \(\vect{AB}\) использовано правило
    \[
    \text{конец}-\text{начало}.
    \]

    \item При вычислении длины отрицательные координаты взяты в скобки:
    \[
    (-3)^2=9.
    \]

    \item Координаты складывались и вычитались по соответствующим компонентам.

    \item При умножении вектора на число умножены обе координаты.

    \item Скалярное произведение записано как
    \[
    x_1x_2+y_1y_2.
    \]

    \item Для косинуса проверено, что оба вектора ненулевые.

    \item Корень не разбит по сумме или разности.

    \item Если требуется квадрат длины, не выполнялись лишние действия с корнем.

    \item Знаки перед отрицательными координатами проверены повторно.
\end{enumerate}

% =================================================
% ТРЕНИРОВОЧНЫЙ БЛОК
% =================================================

\clearpage

\section*{\color{darkaccent}Тренировочный блок --- Путь А}

\sectionline

\textbf{Заданий: 5.}

Решайте последовательно. В тренировочном блоке требуется записать
только решение и итоговый ответ.

\vspace{4mm}

\subsection*{Задание 1}

Даны точки
\[
A(-2;3),
\qquad
B(5;7).
\]

Найдите:

\begin{enumerate}[label=\alph*)]
    \item координаты вектора \(\vect{AB}\);
    \item длину вектора \(\vect{AB}\).
\end{enumerate}

\vspace{7mm}

\subsection*{Задание 2}

Даны векторы
\[
\vec a=(4;-3),
\qquad
\vec b=(-2;5).
\]

Найдите координаты вектора
\[
\vec c=2\vec a-\vec b
\]
и его длину.

\vspace{7mm}

\subsection*{Задание 3}

Даны векторы
\[
\vec p=(3;4),
\qquad
\vec q=(4;-3).
\]

Найдите:

\begin{enumerate}[label=\alph*)]
    \item скалярное произведение \(\vec p\cdot\vec q\);
    \item косинус угла между векторами;
    \item угол между векторами.
\end{enumerate}

\vspace{7mm}

\subsection*{Задание 4}

Даны векторы
\[
\vec a=(2;6),
\qquad
\vec b=(8;4).
\]

Найдите
\[
|\vec a-\vec b|^2.
\]

\vspace{7mm}

\subsection*{Задание 5}

Даны векторы
\[
\vec a=(5;1),
\qquad
\vec b=(2;-3),
\qquad
\vec c=(-1;4).
\]

Найдите значение выражения
\[
(2\vec a-\vec b)\cdot\vec c.
\]

% =================================================
% ОТВЕТЫ
% =================================================

\clearpage

\section*{\color{darkaccent}Ответы}

\sectionline

\vspace{5mm}

\renewcommand{\arraystretch}{1.45}

\begin{table}[ht]
\centering
\begin{tabularx}{\linewidth}{@{}c X@{}}
\toprule
\textbf{№} & \textbf{Ответ} \\
\midrule

1 &
\[
\vect{AB}=(7;4),
\qquad
|\vect{AB}|=\sqrt{65}.
\]
\\

2 &
\[
\vec c=(10;-11),
\qquad
|\vec c|=\sqrt{221}.
\]
\\

3 &
\[
\vec p\cdot\vec q=0,
\qquad
\cos\varphi=0,
\qquad
\varphi=90^\circ.
\]
\\

4 &
\[
|\vec a-\vec b|^2=40.
\]
\\

5 &
\[
(2\vec a-\vec b)\cdot\vec c=-23.
\]
\\

\bottomrule
\end{tabularx}
\end{table}

\vfill

\begin{center}
{\small\color{darkaccent}
Лёвин Артём Александрович
\quad\(\cdot\)\quad
эксклюзивно для Анастасии Павловой
}
\end{center}

\end{document}